% §6 Open problems, and what is elsewhere --- paper-only (Q-0292, 2026-10-04), from the draft's §6
% (draft/spatial-hemigroup-affine.md) and the "Later" list of notes/HANDOFF.md. Each problem is
% stated with the statement it would complete; no statement of the paper depends on an answer.
% Stated as open, never promised (hub RELEASES.md rule 1). Paper VI (spatial-hemigroup-pyramids)
% is unreleased: prose about further work with a TODO(module) comment, never a \ref, and none of
% its results is stated here.
%
% After the blind review (Q-0293; records/affine/REVIEW-paper-sections-1-6.md): the opening restates
% the contribution and points to the trust base instead of certifying the statements; the claim
% that every problem is named in sections 3-5 is dropped (problem 2 is not); problem 4 carries the
% hypothesis "no Gaussian part" of thm:origin-boundedness-d(4) and the logarithmic form at c = d;
% problem 6(c) is put as a question and 6(d) marks its numerical ground; problem 7 is the computed
% index 2 only; the draft's problem 8 (shape selection, singularity theory) is a paragraph of scope,
% not a problem, since no statement of the paper would be completed by it; the third copy of the
% scope disclaimer on the line theory is cut (section 1 "Scope" and rem:isotropic-line keep it).
%
% Not listed, and why: the converses of lem:operators-to-kernels, which were a statement of scope,
% not a problem, and since Q-0314 are lem:kernels-to-operators (section 1, "Scope"). The novelty of
% K_alpha, the budget and the step cone is no longer a gap: Hudson-Mason (1981) and Jurek (1992) were
% acquired and read (Q-0308), and section 4.6's "Prior art" states the claims against them.
%
% Q-0309: "Stated on kernels" dropped (its content is section 1's "Scope"); "Not treated" is the one
% home of singularity theory; problem 7 is the former appendix remark rem:no-covariance (review flag
% 38); the opening narrowed to what the statements give (flag 5) and followed by a short discussion
% of the main theorems (flag 7).
%
% Q-0307 (the split at §3 | §4, adr/0001): problem 4, the origin in dimension d, moved to the
% isotropic paper (paper-iso/04-open-problems.tex) with the theorem it would complete; the opening's
% "carried the isotropic theory of the line to R^d" went with it. The other problems are unchanged.

% Writing-rules pass, group s47 (rule 11; the isotropic paper's section 4 is the model): the
% section is a conclusion, in the order of the questions of section 1, with what it means for shape
% adaptation and the limitations in the section proper (the former "Limits", "Not treated" and "What
% is elsewhere"), and the open problems as its subsection. The file keeps its name.
\section{Conclusion}\label{sec:conclusion}

This paper asked which positive scale spaces indexed by covariance matrices exist, for a covariance
group $G$ and an order $W$, and what each can do for shape adaptation. For the general linear group
the answer runs from the rays, where every rotation-invariant self-decomposable law keeps its affine
orbit (Proposition~\ref{prop:ray-families}), to the Loewner order, where for $d \ge 2$ the Gaussian is
alone and one step of shape-adapted smoothing already forces it (Theorem~\ref{thm:affine-rigidity} and
Proposition~\ref{prop:shape-adapted-apertures}(3)). Between the two a stable law keeps a cascade at a
price in speed of shape change (Theorem~\ref{thm:wedge-axis}(3)), and the orbits whose transform decays
like a power of the frequency, the \Matern{} orbits among them, cascade at most along rays
(Propositions~\ref{prop:log-orbit-rays} and~\ref{prop:matern-steps}). A smaller group admits more
families. With the similarities, non-Gaussian families survive along the Loewner order, with an exponent
% Round 3 (content review), E item 6: the quantity Table 2 bounds, named.
spread of the non-Gaussian part at most four times the eccentricity of the aperture (Propositions~\ref{prop:similarity-stable-family}(3)
and~\ref{prop:similarity-stable-order-kernels}(2)).
% Round-1 external review (Gemini; PLAN-review-response-paper.md, D3): the summary shortened; the
% statements it repeated are in section 1 and Table 2.

\paragraph{What this means for shape adaptation.} The normalization of Lindeberg and G\aa{}rding,
which holds the least smoothing fixed (\sscite{lindeberg1997shape}, \S\S6.1--6.2), is what places the
adapted apertures where the rigidity acts: on the boundary of the Loewner cone of the first stage,
where one step leaves the Gaussian alone. For the Gaussian the adapted record is a further smoothing
of the first stage, whether or not it is computed that way. For every other kernel it is a new
measurement of the image, which the iteration makes anyway, since each round convolves the original
image (\S\ref{sec:landscape-uses}). So, as far as the apertures are concerned, a non-Gaussian kernel
loses nothing the practice uses, and it loses the cascade between its stages. If the cascade is
wanted, the order must shrink, and Theorem~\ref{thm:wedge-axis} prices the shrinking: a stable kernel
keeps a cascade only at a lower speed of shape change (Table~\ref{tab:speed-limit}), and from
$\alpha = 2$ on no non-Gaussian kernel keeps one. The affine orbits whose symbol in the shape's operator is
a decaying rational function without zeros on $[0, \infty)$, the same for every shape, keep none
across shapes, although finitely many elliptic solves in that operator compute them: they cascade at most along rays
(Proposition~\ref{prop:rational-symbols}). A smaller group admits more families, at the price of
% Round 3 (content review), E item 7.
shape selection that transforms with the image, which needs a group that carries every shape to every
other.

The paper has four limitations. Its results concern families of kernels, and through
Lemma~\ref{lem:kernels-to-operators} the operators they define, with exact cascades, and it makes no
statement about images. It treats the apertures a family offers, not how a shape is chosen. In
shape-adapted smoothing the shape at a point is selected through a second-moment matrix at an
integration scale (\S\ref{sec:landscape-uses}), Proposition~\ref{prop:shape-adapted-apertures} reads
only the apertures that this selection produces, and how the selection and its iteration behave for
the non-Gaussian orbits is not treated. It does not treat the fourth use of
\S\ref{sec:landscape-uses}, singularity theory over the larger group. Along shape directions at
fixed volume there is no order (Lemma~\ref{lem:loewner-scale-shape}(2)), so the singularities met
there are bifurcations of a parametrized family; which monotonicity statements hold along the step
set of a family in dimension $d$, and how extrema behave along the shape directions, are not stated
here. And, apart from the measurement of Remark~\ref{rem:matern-grid}, it does not compute the families
on a grid. What it keeps of computation is
Proposition~\ref{prop:rational-symbols}; the computation by recursive filters and in a pyramid is
left to a further paper by the author, on scale-space pyramids, in preparation.
% TODO(module VI): cite spatial-hemigroup-pyramids (Paper VI) when released.
% Round 3 (content review), E i (the author's decision D10): the general linear cell along the stable
% orders is complete (Proposition 4.22), so three cells are complete and three stop short.
Within these limits the three cells of the general linear group in Table~\ref{tab:landscape} are
complete. Along the stable orders the orbit of a profile, with any Gaussian part, is a
family exactly when the lower exponent of the profile is at least the index of the order
(Proposition~\ref{prop:wedge-axis-profiles}). Three cells stop short of a classification. The
similarities along the stable and Loewner orders are described through bounds on the shape of their
kernels. For the diagonal group along the stable orders the condition on the frame's shapes is
necessary, and whether it is sufficient for a family is open (problems~\ref{prob:continuous-fields}
and~\ref{prob:similarity-shape} below). Three more are classified on part of their parameters: the ray
cells of the two smaller groups up to which kernels extend to a continuous field over the shapes, and
the diagonal group with the Loewner order on the frame's shapes (problem~\ref{prob:continuous-fields}).

\subsection{Open problems}\label{sec:open-problems}

Each problem names the statement it would complete.
\begin{enumerate}
% Round 3 (content review), E i (the author's decision D10): the converse inside a family is
% Proposition 4.22; what stays open is the single step. D-R6: "steps" for "cone".
\item \emph{The single-step converse of the ordered-density steps.}\label{prob:ordered-density-converse} Proposition~\ref{prop:orbit-step-cone} gives
a positive increment for every step whose \Levy{} densities are ordered, $n_\Sigma \le n_{\Sigma'}$.
Inside a family along a stable order or a cone order the converse holds for every profile,
since there the increments are infinitely divisible (Proposition~\ref{prop:wedge-axis-profiles}). For a single step, that a
positive increment forces ordered densities is proved for the stable profiles
(Proposition~\ref{prop:stable-step-cone}) and for the profiles bounded at the origin
(Proposition~\ref{prop:bounded-profile-rays}). For the other profiles it is open, the tempered
profiles $r^{-\alpha}j$ of Proposition~\ref{prop:tempered-step-cone} included. The positivity of an
increment does not by itself exclude a signed \Levy{} measure. The converse would make exact the rows
of Table~\ref{tab:speed-limit} that read ``at least''.

\item \emph{The step set of a general orbit.} The step set of an affine orbit is a step set in the
sense of Proposition~\ref{prop:invariant-orders} (\S\ref{sec:affine-stable-cone}). Whether its relation
is a closed partial order whose steps
are closed under the fractional powers $M \mapsto M^\tau$, $0 \le \tau \le 1$, so that by
Theorem~\ref{thm:cone-orders}(2) it is generated by a cone, is open for a general orbit.
% CHANGED Q-0281 wave 2: blind review W2C-1. The obstruction named here, that Levy's continuity theorem
% on R^d is not among the cited facts, is gone: the paper uses that theorem as a standard one (Mathlib,
% the author's decision of 2026-10-05). Closedness is still unproved, and the fractional powers have
% no route, so the problem stays; proving closedness would be a new statement, the author's call.
Closedness asks that a limit of increments be an increment, and \Levy's continuity theorem on
$\RR^d$, in its form for a pointwise limit of transforms continuous at the origin, is the natural
tool for it. The fractional powers have no such route.

\item \emph{Orders that no cone generates.} Theorem~\ref{thm:wedge-axis} describes the general linear
row along the orders of invariant cones. Along a closed invariant order with no cone, of the kind of
Theorem~\ref{thm:cone-orders}(3), the row is not described.
% Round 3 (content review), E i: the second question of this problem is answered by Proposition 4.22(2).
Along the orders of the cones of Theorem~\ref{thm:wedge-axis}(2)(c) the laws $\mu_I$ of the families
depend on the cone only through the number
$\beta(C)$ of Theorem~\ref{thm:wedge-axis}(2)(c) (Proposition~\ref{prop:wedge-axis-profiles}(2)).

\item \emph{Continuous fields over the shapes.}\label{prob:continuous-fields} A family of the similarities with the ray order is
determined by its kernels at the diagonal matrices with nonincreasing nonnegative entries and trace $1$,
one for each orbit of the similarities (Proposition~\ref{prop:similarity-ray-families}(1)--(2)). Which choices of those kernels extend to a
continuous field on $\PSD(d)$ is open. So is the same question over the correlation matrices for the
positive diagonal group (Proposition~\ref{prop:diagonal-ray-families}). Along the stable order the
% Round 3 (content review), E item 10: "the converse" said as what it asks.
families of the diagonal group off the frame's shapes are not described, and whether every law of
Proposition~\ref{prop:diagonal-stable-order-families}(2) is the isotropic member of a family is open. The groups that contain
$\mathrm{SO}(d)$ but not $\Orth(d)$ and are not transitive on $\PD(d)$, such as the
orientation-preserving similarities, are open as well (Appendix~\ref{app:orientation}).

\item \emph{The shape of the similarity families.}\label{prob:similarity-shape} Four questions remain.
(a) For the similarities along the stable order of index $\alpha$, the exponent spread of the
non-Gaussian part of a kernel is at most $4\operatorname{ecc}(\Sigma)^{(d+\alpha)/\alpha}$
(Proposition~\ref{prop:similarity-stable-order-kernels}(2)). The stable orbit has the spread
$\operatorname{ecc}(\Sigma)^{\alpha/2}$, and the sharp exponent is open.
(b) The comparability hypotheses of Proposition~\ref{prop:similarity-general-bounded-shape} are
sufficient and not necessary, and a necessary condition for a bounded shape is open.
(c) Hypothesis (D) of Proposition~\ref{prop:similarity-general-index-zero}(3), for a general part
whose polar profiles are bounded at the origin, asks for the second term of the exponent's
expansion, and it is open. Clause (1) of that proposition gives the index at the origin direction by
direction for infinitely divisible increments. Whether it holds when the increments are only known to
be positive is not known; clause (2) does hold then, when there is no Gaussian part.
(d) For families of kernels that are not cascades, with a Gaussian part and positive increments
only, the infinite divisibility that Proposition~\ref{prop:similarity-orbit-family}(1)--(3) asks
for need not hold: the computation of Remark~\ref{rem:gaussian-matern-step} finds a positive
increment that is not infinitely divisible, a numerical observation. Whether
the shape of such families is bounded is open.

\item \emph{How much the \Matern{} orbit fails.} The mechanism of the failure off the rays, and its
negative mass diverging like a logarithm at the origin, are proved for the \Matern{} orbit of
index $1$ (Remark~\ref{rem:matern-mechanism} and Proposition~\ref{prop:matern-negative-mass}). The
growth that the computation of Remark~\ref{rem:matern-grid} measures on a grid, linear in the
logarithm of the aperture, is not proved on the lattice.
% Round 3 (content review), D-R3 (VERIFY-round3-landscape-orbits-paper.md item 5): the index-2 rate
% is computable; it stays a numerical finding (PLAN-review-response-paper.md R3.1).
At index $2$, the other index computed, the grid follows within 4\,\% the rate that the same
heuristic reads off the degree-$0$ limit of the increment's transform, as at index $1$
(Remark~\ref{rem:matern-grid}). That rate is computed and not proved.

\item \emph{A form of the rigidity without covariance.}\label{prob:no-covariance} The core of
Theorem~\ref{thm:affine-rigidity} that Remark~\ref{rem:kelker} isolates suggests a form without any
covariance group, which is not proved. Its hypotheses are these. A
kernel is reached from scale zero along two decompositions into rank-one steps; the step
directions of each decomposition form a basis of $\RR^d$; the two bases differ; and each increment
is a law carried by the line of its step. Each decomposition writes the kernel as the law of a sum
of independent one-dimensional components along its basis. These are the hypotheses of the
theorem of Darmois and Skitovich and of its matrix form in Comon, with one step missing: the two decompositions
must be read as one vector written twice as a linear image of independent components, and which pairs
of bases count as mixing must be settled. The expected conclusion is that the Darmois--Skitovich
theorem makes Gaussian every non-degenerate component that the two bases mix: in Skitovich's form
(\sscite{skitovich1954linear}, Thm.~2, p.~190, with Thm.~1, p.~185, for the case in which every
component is mixed), and in Darmois's as Comon states it (\sscite{comon1994independent}, Thm.~19,
p.~306), with its matrix form, Thm.~10, p.~294. It would
extend Theorem~\ref{thm:rank-deficient-step} to families with no group.
\end{enumerate}

% The former paragraphs "Not treated" and "What is elsewhere" are the limitations of the section
% proper (group s47).
